\section{Мера Лебега в $\R^n$}

Напомним, что 
\[\mathcal{P}_n := \{[a,b) \ | \ a = (a_1, \dots, a_n), \ b = (b_1, \dots, b_n), \ \forall i \ a_i \leqslant b_i\}\]
— полукольцо.
\[\mu([a,b)) := \prod\limits_{j=1}^{n} \mu([a_j, b_j)) = \prod\limits_{j=1}^{n} (b_j - a_j).\]

\begin{remark}
    Используя рассуждения, аналогичные одномерным, доказывается, что $\mu$ — счётно-аддитивная мера на $\mathcal{P}_n$.
    Кроме того, $\mu$ — $\sigma$-конечная мера на $\mathcal{P}_n$, поскольку
    \[\R^n = \bigcup\limits_{n=1}^{\infty}[-d, d).\]
\end{remark}

Применим всю предыдущую схему к мере $\mu$ на полукольце ячеек и получим
\begin{definition}
    \[\mu^*(E) := \inf \left\{ \sum\limits_{j=1}^{\infty} \mu(P_j) \ \Big| \ \{P_j\} \subset \mathcal{P}, \ E \subset \bigcup\limits_{j=1}^{\infty} P_j\right\}\]
    — \textit{внешняя мера Лебега в $\R^n$}.
     %(далее обозначается $\mathcal{L}_n$}).
\end{definition} 

\begin{definition}
    Минимальная $\sigma$-алгебра, порождённая $\mathcal{P}_n$, называется \textit{борелевской $\sigma$-алгеброй} $\mathcal{B}(\R^n)$.
\end{definition} 

\begin{definition}
    $\mathfrak{M}_n$ — класс всех измеримых множеств в $\R^n$.
\end{definition} 

\noindent \textbf{Аксиома выбора.}
\textit{Пусть $\{E_{\alpha}\}_{\alpha \in I}$ — семейство абстрактных множеств.
Тогда существует функция
\[F: I \to \bigcup\limits_{\alpha \in I} E_{\alpha}: \ F(\alpha) \in E_{\alpha} \ \forall \alpha \in I\]
и
\[E := \{F(\alpha) \ | \ \alpha \in I\}\]
можно считать множеством.}

\begin{theorem}
    Для любого измеримого по Лебегу множества $A \subset \R^n$ положительной меры существует неизмеримое подмножество $E \subset A$.
\end{theorem} 
\begin{proof}
    Введём на множестве $A$ введём отношение эквивалентности:
    \[x \sim y \Longleftrightarrow x - y \in \Q^n.\]
    Получили разбиение $A$ на классы эквивалентности.
    $\{E_{\alpha}\}_{\alpha \in I}$ — семейство классов эквивалентности. 
    Пользуясь аксиомой выбора, построим множество $E \subset A$:
    \[E \cap E_\alpha = \{X_\alpha\} \ \forall \alpha \in I, \ E = \{X_\alpha\}_{\alpha \in I}.\]
    Утверждается, что $E$ неизмеримо.
    Поскольку $A$ ограничено, то
    \[\exists R > 0: \ \forall x, y \in A: \ \|x-y\| \leqslant R.\]
    Рассмотрим всевозможные рациональные сдвиги множества $E$, по модулю не превосходящие $R$.
    По построению
    \[\forall r, r' \in \Q^n: \ (E + r) \cap (E + r') = \emptyset.\]
    С другой стороны,
    \[A \subset \bigcup\limits_{\substack{r \in \Q^n \\ \|r\| \leqslant R}} (E + r).\]
    Предположим, что $E$ измеримо, тогда:
        \[
        \left[
        \begin{aligned}
        \mu(E) = 0 \\
        \mu(E) > 0.
        \end{aligned}
        \right.
    \]
    Случай равенства нулю невозможен, поскольку тогда 
    \[\mu(A) \leqslant \sum\limits_{\|r\| \leqslant R} \mu(E + r) = 0.\]
    Если же $\mu(A) > 0$, то
    \[\mu(E + r) = \mu(E)\]
    — сдвиг не меняет меру.
    Тогда
    \[\mu\left(\bigcup\limits_{\substack{r \in \Q^n \\ \|r\| \leqslant R}} (E + r)\right) = \sum\limits_{n=1}^{\infty} \mu(E + r) = \sum\limits_{n=1}^{\infty} \mu(E) = +\infty.\]
    Однако такого быть не может, поскольку
    \[\bigcup\limits_{\substack{r \in \Q^n \\ \|r\| \leqslant R}} (E + r) \subset B_{2R} (x) \ \forall x \in E,\]
    значит,
    \[\mu\left(\bigcup\limits_{\substack{r \in \Q^n \\ \|r\| \leqslant R}} (E + r)\right) \leqslant \mu(B_{2R}(x)) < + \infty\]
    — противоречие с монотонностью меры.
\end{proof} 

